\section{微扰论中的重正化}\label{sec:theory}

根据局域量子场论中的标准重正化理论~\cite{Collins:1984xc}，算符矩阵元的重正化不依赖于外部态，只与算符本身的短距离性质有关。因此，可以在微扰格林函数或离壳夸克与胶子矩阵元中研究算符的重正化性质。对于 LaMET 在 PDF 中的应用，我们关心的是式~(\ref{eq:operator}) 中的非定域算符 $O_{\Gamma}(z)=\bar{\psi}(0) \Gamma U(0,z) \psi (z)$。这里存在两类发散：与 Wilson 链接（其自能量）相联系的线性发散，以及与涉及 Wilson 线和轻夸克的顶点重正化相联系的对数发散。重正化是乘性的，且只有线性发散带有（线性的）$z$ 依赖。于是重正化后的算符为~\cite{Ji:2017oey,Ishikawa:2017faj,Green:2017xeu, Ji:2020brr}
\bea\label{eq:Reoperator}
O_{\Gamma}(z)_{R}=Z_{O}^{-1} e^{\delta \bar{m} z} O_\Gamma(z),
\eea
其中 $e^{\delta \bar{m} z}$ 包含线性发散，$Z_{O}$ 包含对数发散。除线性发散外 $\delta \bar{m}$ 并非唯一确定，由此引入一种减除方案依赖，它影响着重正化后算符的 $z$ 依赖性~\cite{Ji:2020brr}。

线性发散的质量重正化可以在微扰论中计算。单圈结果与格点作用量无关~\cite{Chen:2016fxx}，
\bea\label{eq:delta_m}
\delta \bar{m} = \frac{2 \pi}{3 a}(\alpha_s+{\cal O}(\alpha_s^2)+...) \ .
\eea
$\alpha_{s}$ 的能标可取 $1/a$ 以与格点结果匹配，
\bea\label{eq:alpha_s}
\alpha_{s}(1/a, \Lambda_{\rm QCD}) = \frac{2 \pi}{b_{0} \ln[1/(a \Lambda_{\rm QCD})]} \ ,
\eea
其中 $b_{0}=11-\frac{2}{3}n_f$（我们取 $n_f$=3）是含 $n_f$ 种费米子时 QCD $\beta$ 函数的单圈值，$\Lambda_{\rm QCD}$ 是非微扰 QCD 标度。计入 $\beta$ 函数的高阶项会得到更复杂的表达式。能标的不同选取相当于计入不同的高阶修正。此外，$\delta \bar{m}$ 中的 $\Lambda_{\rm QCD}$ 与高阶 $\alpha_s$ 项也依赖于格点作用量。

由于红外重整子，微扰论并不收敛，例如参见~\cite{Ji:1995tm,Beneke:1998ui,Bauer:2011ws,Bali:2013pla}。重整子的存在表明 $\delta \bar m$ 中存在一项与 $a$ 无关的非微扰贡献，
\begin{equation}\label{eq:m0}
      \delta\bar m = m_{-1}(a)/a - m_0 \ ,
\end{equation}
其中负号只是约定。求和微扰级数以求得 $m_{-1}(a)$ 时产生的不确定性，与非微扰 $m_0$ 中相同的不确定性相互补偿，使得总和不受重整子不确定性的影响。

$m_0$ 的另一不确定性来自减除方案，或等价地说来自与连续方案的匹配。为了减小减除方案依赖，我们要求重正化后的格点关联函数在短距离 $z$ 处与连续微扰论的 $\overline{\rm MS}$ 结果一致。更具体地说，我们在窗口 $a \ll z < z_S$ 内通过把格点结果匹配到 $\overline{\rm MS}$ 结果来确定 $m_0$，其中 $z_S < 1/\mu$ 是微扰论开始失效的点，$\mu$ 是某个微扰标度。条件 $a \ll z$ 保证格点结果的离散化误差很小。对这个特殊选择的 $m_{0c}$，LaMET 展开在 $1/P^z$ 中不含线性项~\cite{Ji:2020brr}。

单圈水平、维数正规化下，对数发散的重正化因子为~\cite{Ji:2020ect,Constantinou:2017sej}
\bea\label{eq:ZO}
Z_{O} = 1+\frac{3 C_{F} \alpha_{s}}{2 \pi} \frac{1}{4-d},
\eea
其中 $C_{F}=4/3$ 是 $SU(3)$ 基础表示的 Casimir 算符，$d$ 是时空维数。通过求解重正化群方程可以把对数发散重整化（求和）起来，
\bea\label{eq:ZO_RGE}
\frac{d Z_{O}(\epsilon, \mu)}{d \ln (\mu)}=\gamma Z_{O}(\epsilon, \mu),
\eea
其中 $\epsilon=(4-d)/2$，而 $\gamma=-\frac{3 C_{F}}{2 \pi} \alpha_{s}(\mu, \Lambda_{\rm QCD})$ 是领头反常量纲，它与正规化方法无关。式~(\ref{eq:ZO_RGE}) 领头阶解的形式与正规化方案无关；在格点上为
\bea\label{eq:ZO_RS}
Z_{O}(1/a, \mu)=\left(\frac{\ln[1 /(a \Lambda_{\rm QCD})]}{\ln[\mu / \Lambda_{\rm QCD}]}\right)^{\frac{3 C_{F}}{b_0}},
\eea
其中裸紫外（UV）截断为 $1/a$，重正化标度为 $\mu$。

如果在求解重正化群方程式~(\ref{eq:ZO_RGE}) 时进一步考虑次领头对数对反常量纲 $\gamma$ 的贡献，则得到~\cite{Ji:1991pr}
\bea\label{eq:ZO_RS_higher}
Z_{O}(1/a, \mu)=\left(\frac{\ln [1 /(a \Lambda_{\rm QCD})]}{\ln [\mu / \Lambda_{\rm QCD}]}\right)^{\frac{3 C_{F}}{b_0}}\left(1+\frac{d}{\ln[a \Lambda_{\rm QCD}]}\right),\nonumber\\
\eea
其中 $\Lambda_{\rm QCD}$ 与常数 $d$ 依赖于具体的格点作用量。
